\documentclass[11pt,a4paper]{article}

\usepackage[margin=1.8cm]{geometry}
\usepackage{titlesec}
\usepackage{enumitem}
\usepackage{hyperref}
\usepackage{parskip}

\pagenumbering{gobble}

\titleformat{\section}
  {\large\bfseries}
  {}
  {0em}
  {}[\titlerule]

\setlist[itemize]{noitemsep, topsep=0pt}

\begin{document}

% Header
\begin{center}
{\LARGE \textbf{Félix Douaud}} \\[4pt]
Ingénieur logiciel -- C++ / Systèmes \\[6pt]
\href{mailto:felix.douaud@gmail.com}{felix.douaud@epitech.com} \quad | \quad
\href{https://github.com/processssss}{github.com/processssss} \quad | \quad
Phone: +33 6 16 73 37 30
\end{center}

\vspace{6pt}

\section*{Profil}

Étudiant en troisième année à EPITECH, spécialisé en \textbf{C++ et programmation système}.  
Je m’intéresse particulièrement aux \textbf{systèmes temps réel, à l’architecture logicielle, aux moteurs graphiques et aux systèmes audio}.  
Je développe régulièrement des projets personnels autour du rendu graphique et du routing audio, avec une attention particulière portée aux performances et à la modularité.

\section*{Compétences}

\textbf{Langages} \\
C++, C, Python, JavaScript / TypeScript

\textbf{Technologies} \\
Linux, Git, Docker

\textbf{Outils} \\
CMake, Makefile

\textbf{Domaines} \\
Programmation système, architecture logicielle, moteurs graphiques, programmation audio

\section*{Projets}

\textbf{Kablr -- Moteur de routing audio}

\begin{itemize}
\item Développement d’un moteur de routing audio modulaire
\item Architecture multiplateforme (Linux / Windows)
\item Implémenté en C++20 avec RtAudio
\item Routage dynamique entre périphériques audio
\end{itemize}

\textbf{Raytracer -- Moteur de rendu}

\begin{itemize}
\item Implémentation d’un moteur de rendu basé sur le ray tracing en C++
\item Parsing de scènes et simulation de l’éclairage
\item Architecture orientée objet pour le moteur de rendu
\end{itemize}

\textbf{AREA -- Plateforme d’automatisation}

\begin{itemize}
\item Plateforme d’automatisation inspirée de IFTTT
\item Architecture backend en Node.js
\item Déploiement conteneurisé avec Docker
\end{itemize}

\section*{Formation}

\textbf{EPITECH Nancy -- 3\textsuperscript{e} année} \\
Diplôme d’Expert en Ingénierie Logicielle \hfill 2023 -- 2028

\begin{minipage}[t]{0.48\textwidth}

\section*{Langues}

Français -- Natif \\
Anglais -- B1

\end{minipage}
\hfill
\begin{minipage}[t]{0.48\textwidth}

\section*{Centres d’intérêt}

Judo -- Ceinture noire 1\textsuperscript{er} dan \\
Systèmes Linux \\
Jeux vidéo compétitifs

\end{minipage}

\end{document}
